\problemname{The Powers of Two}

Louise and her close friend Elin have always been completely unstoppable together.
They called their unbeatable duo-energy the Power of Friendship, and when they worked together, nothing could stop them.

But one day, Louise discovered something even greater. Something so big that it might not fit in a 32-bit integer.

She was counting by the campfire and noticed that certain numbers grew in a special way: $1, 2, 4, 8, 16, 32, \dots$  
Each number was twice as large as the previous one, which makes the numbers grow very quickly.

\begin{align*}
& \text{``This... this is more than the power of friendship'', whispered Louise.} \\
& \text{``It is the power of powers of two''.}
\end{align*}

Yes. The power of powers of two is truly something special. But now Louise and Elin wonder whether a given number $N$ will ever appear in the sequence of powers of two.

\section*{Input}
The input consists of a single integer $N$ ($1 \leq N \leq 10^{18}$).

\section*{Output}
Print \verb|Yes| if $N$ is a power of two, otherwise print \verb|No|.

\section*{Scoring}
Your solution will be tested on a set of test groups.
To earn points for a group, you must pass all test cases in that group.

\noindent
\begin{tabular}{| l | l | p{12cm} |}
  \hline
  \textbf{Group} & \textbf{Points} & \textbf{Constraints} \\ \hline
  $1$    & $20$  & $N \le 100$ \\ \hline
  $2$    & $80$  & No additional constraints. \\ \hline
\end{tabular}
